\section*{Related Work and Conceptual Framing}

Prism adaptation is usually described through two linked behavioural signatures. \emph{Error reduction} refers to the reduction of pointing error during exposure to the displacement, while the \emph{after-effect} refers to a remaining shift in behaviour once the displacement is removed \cite{serino2006mechanisms,Serino2011}. Open-loop pointing is often used to measure sensorimotor after-effects because participants point without the same online visual feedback used during exposure \cite{Rossetti1998Prism,Ramos2019,anan2025comparative}. Line bisection and landmark tasks are used to inspect whether adaptation transfers to spatial judgement rather than only changing endpoint motor behaviour \cite{frassinetti2002long,Serino2011,Bourgeois2021}.

The procedure around the perturbation also matters. Ladavas et al. compared terminal prism adaptation, where only the final part of the movement is visible, with concurrent prism adaptation, where more online feedback is available during movement. Both procedures improved neglect relative to a control, but terminal prism adaptation produced stronger clinical improvement \cite{Serino2011}. Carter et al. similarly hid the initial part of reaching in a virtual-shift procedure because visual feedback during the movement can reduce adaptation induction \cite{Carter2016}. Bourgeois et al. further showed that visual feedback, but not auditory-verbal feedback, induced after-effects in a virtual prism procedure \cite{Bourgeois2021}. These findings show that feedback visibility and modality are not minor presentation details; they are part of the adaptation procedure.

VR studies vary in how the prism-like mismatch is implemented. Kim et al. and Cho et al. used virtual hand displacement to reproduce key features of prism adaptation without physical prisms \cite{kim2017development,cho2022feasibility}. Ramos et al. compared physical prisms with two immersive VR manipulations and reported larger after-effects in VR conditions than in the physical-prism condition \cite{Ramos2019}. Anan et al. compared conventional prism glasses with a VR prism-adaptation system and reported that VR induced leftward open-loop pointing shifts in more participants \cite{anan2025comparative}. Gerken et al. studied visuomotor rotation and found that hand-based feedback could speed early error reduction without necessarily changing after-effects or transfer \cite{gerken2025sense}. W{\"a}hnert and Sch{\"a}fer further showed that participant instructions can influence after-effect persistence in VR, reinforcing that procedure and expectation are part of the experimental setup \cite{Wahnert2024}.

The related work therefore contains two separable design questions. The first is how the visual mismatch is created: some systems shift a scene or rendering reference, some displace the virtual effector or response endpoint, and others rotate the relation between movement and feedback. Appendix Table~\ref{tab:article-appendix-effect-categories} summarizes these perturbation categories as they are used to organize the reviewed VR implementations. The second question is how adaptation is measured after exposure. Appendix Table~\ref{tab:article-appendix-task-categories} summarizes the assessment tasks that recur in the reviewed literature and in the prototype design.

Based on this diversity, the present project separates three implementation categories. \emph{Translation} shifts the response origin laterally. \emph{Rotation} changes the direction of the pointing response. \emph{Skew} shifts the visible pointing representation together with its corresponding aim, making the virtual hand or controller itself appear displaced. These labels are prototype-level implementation labels rather than claims that the three manipulations are physiologically equivalent. The category labels and the appendix tables are therefore used as synthesis devices: they make explicit which kind of perturbation is being implemented and which behavioural task is used to inspect the resulting after-effect or transfer.
